\chapter{TÍCH HỢP VÀO NSGA-II VÀ KIỂM THỬ}
\label{chap:tichhop}

\section{Vị trí gắn trong vòng lặp}

Hình~\ref{fig:flow} là luồng một lần chạy NSGA-II sau khi gắn hai cải tiến. Ô nét đứt là phần chỉ chạy
khi cờ tương ứng bật; khi tắt, luồng đi đúng các ô của thuật toán gốc.

\begin{figure}[H]
\centering
\begin{tikzpicture}[node distance=0.55cm and 0.9cm,
  stp/.style={draw, rounded corners=2pt, minimum width=5.2cm, minimum height=0.75cm,
               font=\small, align=center},
  tv2/.style={stp, dashed, draw=hustred, fill=hustred!6},
  tv4/.style={stp, dashed, draw=grayline, fill=black!4},
  dec/.style={draw, diamond, aspect=2.6, inner sep=1pt, font=\small, align=center},
  arr/.style={-{Stealth[length=2mm]}, thick}]
  \node[dec] (c1) {\code{init\_smart} bật?};
  \node[tv2, minimum width=5.6cm] (i1) at ($(c1)+(-3.4,-1.6)$) {quần thể NN / TW / ngẫu nhiên\\ (\code{smart\_initial\_population})};
  \node[stp, minimum width=5.6cm] (i0) at ($(c1)+(3.4,-1.6)$) {$N_{pop}$ hoán vị ngẫu nhiên\\ (thuật toán gốc)};
  \node[stp] (ev0) at ($(c1)+(0,-3.1)$) {đánh giá quần thể đầu};
  \node[stp, below=of ev0] (sel) {giải đấu nhị phân chọn cha mẹ};
  \node[stp, below=of sel] (cx) {lai ghép OX (xác suất $p_c$)};
  \node[stp, below=of cx] (mu) {đột biến (xác suất $p_m$)};
  \node[tv4, below=of mu] (w) {gán trọng số DP cho con (TV4, nếu bật)};
  \node[tv2, below=of w] (ls) {với xác suất $p_{ls}$: tìm kiếm cục bộ\\ (\code{local\_search\_children})};
  \node[stp, below=of ls] (ev) {đánh giá các con chưa có kết quả};
  \node[stp, below=of ev] (sur) {gộp cha mẹ và con, chọn lọc môi trường};
  \draw[arr] (c1.west) -| node[pos=0.25, above, font=\footnotesize] {có} (i1.north);
  \draw[arr] (c1.east) -| node[pos=0.25, above, font=\footnotesize] {không} (i0.north);
  \draw[arr] (i1.south) |- (ev0.west);
  \draw[arr] (i0.south) |- (ev0.east);
  \foreach \a/\b in {ev0/sel, sel/cx, cx/mu, mu/w, w/ls, ls/ev, ev/sur} \draw[arr] (\a) -- (\b);
  \draw[arr] (sur.east) -- ++(1.1,0) |- node[pos=0.25, right, font=\footnotesize, align=left]
        {lặp $G$\\ thế hệ} (sel.east);
\end{tikzpicture}
\caption{Luồng NSGA-II của nhóm và hai điểm gắn của TV2 (viền đỏ nét đứt)}
\label{fig:flow}
\end{figure}

\section{Điểm gắn khởi tạo}

Trong \code{nsga2.run}, dòng tạo quần thể đầu của thuật toán gốc được bọc trong một nhánh điều kiện:
\begin{Verbatim}
if imp.get("init_smart", False):   # --- TV2 init_smart
    ids = np.arange(1, inst.n + 1)
    perms, init_sources = smart_initial_population(
        imp, inst, n_pop, rng, lambda r: r.permutation(ids).tolist())
    pop = [np.array(x, dtype=ids.dtype) for x in perms]
else:
    pop = [rng.permutation(np.arange(1, inst.n + 1)) for _ in range(n_pop)]
\end{Verbatim}
Hàm sinh cá thể ngẫu nhiên truyền vào chính là dòng sinh của thuật toán gốc, nên nhóm \code{random}
trong quần thể thông minh có cùng phân phối với thuật toán gốc. Nhánh \code{else} giữ nguyên văn dòng
cũ.

\section{Điểm gắn tìm kiếm cục bộ}

Đối tượng \code{LocalSearchHook} chỉ được tạo khi \code{improvements.local\_search} bật. Bộ đánh giá
truyền cho nó là đúng bộ đánh giá NSGA-II đang dùng: \code{evaluate\_detail} khi không bật DP,
\code{evaluate\_dp} với trọng số của con khi bật DP. Sau khi \code{make\_offspring} trả về các con (đã
đột biến và đã có trọng số DP nếu bật), hàm \code{local\_search\_children} xử lý lần lượt từng con:
\begin{enumerate}
  \item gọi hook: bốc một số ngẫu nhiên, nếu nhỏ hơn $p_{ls}$ thì chạy tìm kiếm cục bộ;
  \item con đã qua tìm kiếm cục bộ được thay bằng hoán vị kết quả, và vectơ mục tiêu lấy luôn từ kết quả
        đánh giá chính xác mà tìm kiếm cục bộ trả về;
  \item các con còn lại được đánh giá đúng như thuật toán gốc.
\end{enumerate}

\paragraph{Đếm số lần đánh giá.} Lần chấm nghiệm ban đầu bên trong tìm kiếm cục bộ thay cho lần chấm
chính thức của con, nên con đó không bị chấm hai lần. Số lần đánh giá của một lần chạy là
\begin{equation}
  n_{\text{evals}} = N_{pop}\,(G + 1) + R - T,
  \label{eq:nevals}
\end{equation}
với $R$ là tổng số lần tìm kiếm cục bộ thật sự gọi bộ đánh giá và $T$ là số lần tìm kiếm cục bộ được
kích hoạt. Công thức này được kiểm tra trực tiếp trong test tích hợp.

\section{Quy tắc số ngẫu nhiên và nhánh cờ tắt}

Hai cải tiến chỉ dùng bộ sinh \code{rng} của \code{nsga2.run}. Khi cờ tắt:
\begin{itemize}
  \item \code{init\_smart}: không gọi mô-đun, không đọc tham số, dòng tạo quần thể giữ nguyên;
  \item \code{local\_search}: không tạo hook, không có lần bốc $p_{ls}$ nào, các con được đánh giá bằng
        đúng lời gọi \code{evaluate\_population} như trước.
\end{itemize}
Để tái dùng kết quả đánh giá của tìm kiếm cục bộ, hàm chuyển một kết quả bộ giải mã thành vectơ mục tiêu
và mức vi phạm số xe được tách ra thành \code{objective\_row} và dùng chung cho cả \code{evaluate\_population}.
Phép tính bên trong giữ nguyên, nên kết quả của thuật toán gốc không đổi đến từng bit (mục~\ref{sec:regress}).

\section{Tương thích với các cải tiến của TV3 và TV4}

\begin{itemize}
  \item \textbf{TV3 (\code{diversity})}: tác động ở bước chọn lọc và kho lưu trữ, sau bước đánh giá con.
        Tìm kiếm cục bộ trả về vectơ mục tiêu của con trước khi kho lưu trữ được cập nhật, nên hai cải tiến
        không chạm vào nhau.
  \item \textbf{TV4 (\code{dp\_split} chế độ \code{every\_eval})}: tìm kiếm cục bộ chấm bằng
        \code{evaluate\_dp} với đúng trọng số của con, không ép bộ tuyến, không bốc lại trọng số.
  \item \textbf{TV4 (\code{dp\_split} chế độ \code{polish})}: quá trình tiến hoá dùng cách cắt gốc, DP chỉ
        chạy sau thế hệ cuối, nên tìm kiếm cục bộ vẫn kiểm tra lệch biên như khi không bật DP.
  \item \textbf{TV4 (\code{adaptive\_rates})}: chỉ đổi $p_c$, $p_m$ theo từng cặp cha mẹ, không liên quan
        tới tìm kiếm cục bộ.
\end{itemize}
Các tổ hợp này đều có test chạy NSGA-II thật (bảng~\ref{tab:tests}).

\section{Kiểm thử}

\subsection{Bộ test}

Toàn bộ dự án có 432 test chạy bằng \code{pytest}, trong đó 206 test thuộc phần việc TV2.

\begin{table}[H]
\centering
\caption{Các nhóm test của TV2}
\label{tab:tests}
\begin{tabularx}{\textwidth}{@{}lrX@{}}
\toprule
\textbf{Tệp} & \textbf{Số test} & \textbf{Nội dung kiểm tra} \\
\midrule
\code{test\_init\_smart.py} & 91 &
  mọi cá thể là hoán vị đúng; chia số lượng 16/16/48 và các $N_{pop}$ lẻ; phá hoà NN và TW trên dữ liệu dựng
  tay; sắp theo $ETa$ chứ không theo $ET$; phép hoán đổi giữ tập khách và gọi số ngẫu nhiên đúng số lần;
  $n = 1..5$ không treo; nhãn \code{random\_fallback}; cùng seed ra cùng quần thể; không sửa đầu vào;
  tham số sai báo lỗi; nhóm ngẫu nhiên trùng khớp dãy của thuật toán gốc \\
\code{test\_local\_search.py} & 79 &
  2-opt và Or-opt theo ví dụ và các biên; bước bốc luôn hợp lệ, không vô tác dụng; trội có $\varepsilon$;
  $f_3$ đúng chiều cực đại; kết quả đúng bằng chấm lại độc lập; mỗi bước nhận đều làm giảm điểm hoặc trội;
  mỗi ứng viên chỉ đổi đúng một tuyến và giữ tập khách; phản ví dụ lệch biên bị loại; trọng số bốc
  đúng một lần; truyền nguyên \code{metadata}; trần ngân sách; bộ nhớ đệm; chính sách số xe; bộ đánh giá
  kiểu DP; hook với $p_{ls} = 0$ chỉ bốc một số và không chấm thêm \\
\code{test\_tv2\_integration.py} & 34 &
  trên NSGA-II thật: cờ tắt không gọi code TV2; các tổ hợp \code{init\_smart}, \code{local\_search},
  \code{diversity}, \code{dp\_split}, \code{adaptive\_rates} chạy được, tái lập theo seed, vectơ mục tiêu
  khớp chấm lại độc lập, số lần đánh giá đúng công thức~\ref{eq:nevals}; mốc hồi quy; các hàm phụ của
  harness \\
\code{test\_plot\_tv2.py} & 2 & script vẽ biểu đồ chạy được và báo lỗi khi thiếu cột dữ liệu \\
\bottomrule
\end{tabularx}
\end{table}

\subsection{Kiểm tra hồi quy khi tắt cờ}
\label{sec:regress}

Điều kiện quan trọng nhất của mọi cải tiến là khi tắt hết cờ, kết quả phải giống thuật toán gốc không
chỉ gần đúng mà trùng từng bit. Có hai mốc được kiểm tra:
\begin{itemize}
  \item \code{tests/fixtures/tv4\_alloff\_main.json} do TV4 sinh từ \code{nsga2.py} chưa sửa của nhánh
        chính: hoán vị cuối, vectơ mục tiêu lưu dạng \code{float.hex}, mức vi phạm, front và số lần đánh giá
        trên C102-25 và R201-50, seed 0 và 1. Sau khi gắn hai cải tiến của TV2, cả bốn trường hợp vẫn trùng
        khớp hoàn toàn.
  \item \code{tests/fixtures/tv2\_alloff\_snapshot.json}: mốc của TV2 trên ba instance, hai seed, gồm cả
        trạng thái bộ sinh số ngẫu nhiên cuối cùng.
\end{itemize}
Ngoài ra có test thay các hàm của TV2 bằng hàm ném lỗi rồi chạy NSGA-II với cờ tắt; test pass chứng tỏ
nhánh cờ tắt không chạm tới mã của TV2.

\subsection{Kết quả chạy kiểm thử}

\begin{table}[H]
\centering
\caption{Kết quả các lệnh kiểm thử trên nhánh \code{feat/tv2-init-local-search} sau khi ghép nhánh chính}
\label{tab:testrun}
\begin{tabularx}{\textwidth}{@{}Xl@{}}
\toprule
\textbf{Lệnh} & \textbf{Kết quả} \\
\midrule
\code{python -m pytest} & 432 passed \\
\code{python -m scripts.tv2\_regression\_snapshot --check} & REGRESSION PASS \\
\code{python -m scripts.tv2\_dp\_compat\_check} & DP+LS compatibility: PASS \\
\code{python -m scripts.smoke\_phase0} & chạy thành công \\
\code{python -m pip check} & No broken requirements found \\
\bottomrule
\end{tabularx}
\end{table}

Script \code{tv2\_dp\_compat\_check} chạy tìm kiếm cục bộ với bộ đánh giá DP của TV4 trên C102-25,
R201-50, RC201-100, mỗi instance 20 lần gọi với $p_{ls} = 1$. Ở mọi lần gọi, bộ đánh giá nhận đúng một
đối tượng trọng số, kết quả trả về trùng với chấm lại bằng \code{evaluate\_dp}, và không lần nào biến
nghiệm hợp lệ về số xe thành không hợp lệ. Thời gian trung bình mỗi lần gọi là 21 ms (25 khách),
54 ms (50 khách) và 124 ms (100 khách).
